\section{gpt pro}
\section{Nov 7, try 1}
\section{On the class \texorpdfstring{$\Nu_{\text{good}}$}{Nu\_good} and a parametric family \texorpdfstring{$\Nu(\Theta)$}{Nu(Theta)}}

Throughout, fix a Legendre potential $F$ with domain $\cD\subseteq\R^d$, learning rate $\eta>0$, and a (possibly data–dependent) point $a\in \cA\cap\interior(\cD)$.

\subsection{Generative model and posteriors}
We take the \emph{prior} to be a probability measure $\nu\in\cP(\cA\times\cL)$ over pairs $(A_\star,\ell)$ (optimal action and loss vector). A sampling rule $P\in\Delta(\cA)$ produces $A\sim P$, conditionally independent of $(A_\star,\ell)$ under $\nu$. The observation is
\[
    Y\;=\;\ip{A,\ell}\in\R.
\]
This induces a joint law $\nu^P$ on $(A_\star,\ell,A,Y)$ via
\[
    (A_\star,\ell)\sim \nu,\qquad A\sim P\ \perp\ (A_\star,\ell),\qquad Y=\ip{A,\ell}.
\]
All conditionals such as $\E[\cdot\mid A,Y]$ are with respect to $\nu^P$.
We use the shorthand
\[
    \pr \; := \; \E[A_\star],\qquad \po \;:=\; \E[A_\star\mid A,Y].
\]

\subsection{Multiplicative-\texorpdfstring{$Y$}{Y} estimators and the mean-path map}
Given $T:\cA\to\R^d$, the multiplicative-$Y$ estimator is $\hl=Y\,T_A$ and the induced mirror update is
\[
    \tilde A \;:=\; {F^\star}'\!\big(F'(a)-\eta\,\hl\big)\;=\;{F^\star}'\!\big(F'(a)-\eta\,Y\,T_A\big).
\]
We are interested in when multiplicative-$Y$ can match the unrestricted optimum (i.e., attain the same inner value in the biased–SMD bound). The critical identity is $\tilde A=\po$ almost surely.

\begin{definition}[The good class $\Nu_{\text{good}}$]
    \label{def:Nu-good}
    We say $\nu\in\Nu_{\text{good}}$ if there exists a measurable $T:\cA\to\R^d$ such that
    \[
        {F^\star}'\!\big(F'(a)-\eta\,Y\,T_A\big)\;=\;\E[A_\star\mid A,Y]\qquad\text{a.s. under }\nu^P.
    \]
    Equivalently (apply $F'$ to both sides), for $\nu^P$-a.e.\ $A$ there exists $T_A$ with
    \begin{equation}
        \label{eq:mean-path}
        F'\!\big(\E[A_\star\mid A, Y=y]\big)\;=\;F'(a)\;-\;\eta\,y\,T_A\qquad\text{for all $y$ in the support of $Y\!\mid\!A$.}
    \end{equation}
    We call \eqref{eq:mean-path} the \emph{mirror–affine mean-path condition}.
\end{definition}

\begin{proposition}[Exact characterization of $\Nu_{\text{good}}$]
    \label{prop:Nu-good-iff}
    For Legendre $F$, $\nu\in\Nu_{\text{good}}$ if and only if the mirror–affine mean-path condition \eqref{eq:mean-path} holds (for $\nu^P$-a.e.\ $A$). In particular, $\Nu_{\text{good}}$ is precisely the set of priors whose conditional posterior mean, mapped through $F'$, is affine in $y$ along each action $A$.
\end{proposition}
\begin{proof}
    If $\tilde A=\po$ a.s., then $F'(\po)=F'(a)-\eta Y T_A$, so $y\mapsto F'(\E[A_\star\mid A,y])$ is affine with slope $-\eta T_A$, proving \eqref{eq:mean-path}. Conversely, if \eqref{eq:mean-path} holds, then $\E[A_\star\mid A,Y]= (F^\star)'(F'(a)-\eta Y T_A)=\tilde A$ a.s.
\end{proof}

\subsection{A parametric family \texorpdfstring{$\Nu(\Theta)$}{Nu(Theta)} with log-partition \texorpdfstring{$F^\star$}{F*}}
We now introduce a parametric family that \emph{guarantees} the equality $\tilde A=\po$ and uses $F^\star$ as the log-partition.

\begin{definition}[Mirror–exponential tilt family]
    \label{def:mirror-exponential}
    Let $\mu$ be a base $\sigma$-finite measure supported on $\cD$. We say $\nu\in\Nu(\Theta)$ if there exists a measurable $T:\cA\to\R^d$ such that the conditional law of $A_\star$ given $(A=a,Y=y)$ has $\mu$-density
    \begin{equation}
        \label{eq:mirror-exp-density}
        x\ \mapsto\ \exp\!\Big( \ip{F'(a)-\eta\,y\,T_A,\ x}\;-\;F^\star(F'(a)-\eta\,y\,T_A)\Big).
    \end{equation}
\end{definition}

\begin{proposition}[$\Nu(\Theta)\subseteq \Nu_{\text{good}}$]
    \label{prop:NuTheta-in-Nugood}
    If $\nu\in\Nu(\Theta)$, then for $\nu^P$-a.e.\ $(A,Y)$,
    \[
        \po \;=\; \E[A_\star\mid A,Y] \;=\; (F^\star)'\!\big(F'(a)-\eta Y T_A\big) \;=\; \tilde A,
    \]
    hence $\nu\in\Nu_{\text{good}}$.
\end{proposition}
\begin{proof}
    In the exponential family with natural parameter $\theta=F'(a)-\eta y T_A$ and log-partition $F^\star(\theta)$, the expectation map is $(F^\star)'(\theta)$, which equals $\tilde A$.
\end{proof}

\begin{remark}[On the role of the sampling rule]
    The inclusion $\Nu(\Theta)\subseteq \Nu_{\text{good}}$ is a statement \emph{after} we augment $\nu$ with the sampling kernel $A\sim P$ and the observation map $Y=\ip{A,\ell}$. We do not alter the prior $\nu$ on $(A_\star,\ell)$; rather, we require that the resulting \emph{conditional} $A_\star\mid(A,Y)$ lie in the mirror–exponential tilt family \eqref{eq:mirror-exp-density} for some measurable $T$.
\end{remark}

\paragraph{When is \texorpdfstring{$\Nu(\Theta)$}{Nu(Theta)} nonempty?}
We need $F^\star$ to be a valid log–Laplace transform relative to a base measure $\mu$ on $\cD$. This holds in many standard cases.

\subsubsection*{Specialization I: Quadratic/Gaussian}
Let $F(x)=\tfrac12 x^\top Q x$ with $Q\succ0$. Then $F'(x)=Qx$, $F^\star(\theta)=\tfrac12\theta^\top Q^{-1}\theta$, $(F^\star)'(\theta)=Q^{-1}\theta$. Choosing $\mu$ as the Gaussian base (Lebesgue with quadratic log–density) makes \eqref{eq:mirror-exp-density} a Gaussian:
\[
    A_\star \mid (A=a,Y=y)\ \sim\ \cN\!\Big(Q^{-1}\!\big(F'(a)-\eta y T_A\big),\; Q^{-1}\Big).
\]
Hence
\(
\po \;=\; Q^{-1}(F'(a)-\eta Y T_A)\;=\; a - \eta Y\,Q^{-1}T_A\;=\;\tilde A.
\)

\subsubsection*{Specialization II: Simplex/softmax}
Let $\cD=\Delta(d)$ and $F(p)=\sum_i p_i\log p_i$ (negentropy). Then
\[
    F^\star(\theta)=\log\sum_{i=1}^d e^{\theta_i},\qquad (F^\star)'(\theta)=\operatorname{softmax}(\theta).
\]
Let $\mu$ be counting measure on $\{e_1,\dots,e_d\}$. Then \eqref{eq:mirror-exp-density} yields the categorical posterior
\[
    A_\star\mid(A=a,Y=y)\ \sim\ \text{Categorical}\!\Big(\operatorname{softmax}\big(\log a - \eta y T_A\big)\Big),
\]
with posterior mean
\(
\po=\operatorname{softmax}(\log a - \eta Y T_A)=\tilde A.
\)
Equivalently, the log-odds satisfy
\[
    \log\frac{\po(i)}{\po(j)} \;=\; \log\frac{a_i}{a_j}\;-\;\eta\,Y\,(T_{A,i}-T_{A,j})\qquad(\forall i,j),
\]
which is the mirror–affine mean-path condition \eqref{eq:mean-path}.

\subsection{Is \texorpdfstring{$\Nu(\Theta)=\Nu_{\text{good}}$}{Nu(Theta)=Nu\_good}?}
In general, $\Nu(\Theta)\subsetneq \Nu_{\text{good}}$. The definition of $\Nu_{\text{good}}$ constrains only the \emph{posterior mean path} $y\mapsto \E[A_\star\mid A,y]$ through \eqref{eq:mean-path}; many non–exponential conditional laws share the same mean path and hence belong to $\Nu_{\text{good}}$ while not lying in the exponential family \eqref{eq:mirror-exp-density}.

\subsection{Max-entropy justification for \texorpdfstring{$\Nu(\Theta)$}{Nu(Theta)}}
Among all conditional distributions of $A_\star$ given $(A=a,Y=y)$ whose mean matches a prescribed mirror–affine path,
\[
    \E[A_\star\mid A=a,Y=y]\;=\;(F^\star)'\!\big(F'(a)-\eta y T_A\big),
\]
the mirror–exponential tilt \eqref{eq:mirror-exp-density} uniquely maximizes the (conditional) entropy relative to $\mu$:
\[
    \text{maximize}\quad -\!\int \log\frac{\mathrm{d}p}{\mathrm{d}\mu}\,\mathrm{d}p
    \quad\text{subject to}\quad
    \int x\,\mathrm{d}p(x)=(F^\star)'\!\big(F'(a)-\eta y T_A\big).
\]
This is the standard maximum–entropy with linear moment constraints; its Lagrange dual uses the convex conjugate of the entropy, i.e., the log–Laplace $F^\star$, and the optimizer is exactly the exponential form \eqref{eq:mirror-exp-density}. Thus $\Nu(\Theta)$ is the canonical ``least–informative'' (maximum–entropy) element of $\Nu_{\text{good}}$ consistent with the required mean path.

\subsection{Practical membership tests and recovery of \texorpdfstring{$T$}{T}}
Given $\nu\in\cP(\cA\times\cL)$ and a sampling rule $P$:
\begin{enumerate}
    \item Compute $m_A(y):= \E[A_\star\mid A, Y=y]$ for $\nu^P$-a.e.\ $A$.
    \item Test mirror–affinity: verify that for the $y$-support of $Y\mid A$,
          \[
              F'\!\big(m_A(y)\big)=F'(a)-\eta y\,T_A
              \quad\text{for some $T_A$ independent of $y$.}
          \]
    \item If true, then $\nu\in\Nu_{\text{good}}$ and any $T$ satisfying
          \(
          T_A = \frac{F'(a)-F'(m_A(y))}{\eta y}
          \)
          (for any $y\neq 0$) is valid; plugging this $T$ into the multiplicative-$Y$ estimator yields $\tilde A=\po$ a.s.
\end{enumerate}
In the simplex/softmax case this reduces to checking that all pairwise log–odds $\log (m_A(y)_i/m_A(y)_j)$ are affine in $y$ with a common slope vector across $(i,j)$.


\clearpage
\section{Nov 5, try 1}
\subsection{Restricting the adversary to a dual-affine class without changing the dual value}

We formalize a class of priors under which the multiplicative-$Y$ estimators attain the
same value as the unrestricted optimal estimators, and show that \emph{on the dual (inf--sup) side},
restricting the adversary to this class incurs \emph{no loss}.

\begin{definition}[Dual-affine posterior classes]\label{def:Nu-aff}
    Fix $a\in\cA$ and $F$ Legendre. For an action $b\in\cA$, let $Y_b \triangleq \ip{\ell,b}$.

    \begin{enumerate}
        \item The \emph{dual-affine posterior class} is
              \[
                  \cN_{\mathrm{aff}}
                  \;\triangleq\;
                  \Big\{
                  \nu\in\cP(\cA\times\cL)\,:\,\exists\,U\in\R^d,\;\{S_b\in\R^d\}_{b\in\cA}
                  \text{ s.t. }F'\!\big(\E[A_\star\mid A=b,Y_b=y]\big)=U+S_b\,y
                  \ \text{a.s. for all $b$ with $P(A=b)>0$}
                  \Big\}.
              \]
        \item The \emph{degenerate dual-affine class} is
              \(
              \cN_{\mathrm{deg}}\triangleq\{\delta_{(a_\star,\ell)}:\ (a_\star,\ell)\in\cA\times\cL\}.
              \)
    \end{enumerate}
\end{definition}

\begin{remark}
    $\cN_{\mathrm{aff}}$ is tailored to multiplicative-$Y$ because, by
    Prop.~\ref{prop:multY-iff}, $\tilde{A}=\po$ a.s.\ iff $F'(\po)=F'(a)-\eta\,Y\,T_A$.
    Thus $U=F'(a)$ and $S_b=-\eta\,T_b$ make multiplicative-$Y$ achieve $\tilde{A}=\po$.
    Also, every Dirac prior is in $\cN_{\mathrm{aff}}$ (the conditional support of $Y_b$ is a singleton, hence the map in Definition~\ref{def:Nu-aff} is trivially affine with $S_b=0$ and $U=F'(a_\star)$), so $\cN_{\mathrm{deg}}\subset\cN_{\mathrm{aff}}$.
\end{remark}

We next isolate the structural property that lets us \emph{restrict the dual adversary to $\cN_{\mathrm{aff}}$ without changing the value}.

\begin{assumption}[Compactness and continuity]\label{ass:compact}
    $\cA$ and $\cL$ are compact, $F$ is Legendre, and $\eta>0$ is fixed.
    For multiplicative estimators we choose $T$ from a convex compact set $\cT\subset(\R^d)^{\cA}$,
    which is w.l.o.g.\ under mild coercivity (the Bregman penalty makes the objective inf-compact in $T$).
\end{assumption}

\begin{lemma}[Linearity of the payoff in $\nu$]\label{lem:linear-in-nu}
    Fix $a,\eta$, $P\in\Delta(\cA)$ and $T\in\cT$. Then
    $\nu\mapsto \Lambda_{a,\eta}(P,E_T\Vert \nu)$ is an affine and continuous functional on $\cP(\cA\times\cL)$ (under weak convergence).
\end{lemma}

\begin{proof}
    Use the alternative form \eqref{eq:Lambda-simplified}:
    \(
    \Lambda = \tfrac{1}{\eta}\big(\Delta(\nu,P) + \E_\nu[\langle A_\star,\hl\rangle + \eta^{-1}F^\star(F'(a)-\eta\hl) - \eta^{-1}F^\star(F'(a))]\big),
    \)
    where the dependence on $\nu$ is only through expectations of bounded continuous functions by Assumption~\ref{ass:compact}.
\end{proof}

\begin{proposition}[Pointwise equality of adversarial best responses]\label{prop:sup-over-aff}
    For any fixed $a,\eta$, $P\in\Delta(\cA)$ and $T\in\cT$,
    \[
        \sup_{\nu\in\cP(\cA\times\cL)} \Lambda_{a,\eta}(P,E_T\Vert \nu)
        \;=\;
        \sup_{\nu\in\cN_{\mathrm{aff}}} \Lambda_{a,\eta}(P,E_T\Vert \nu).
    \]
\end{proposition}

\begin{proof}
    By Lemma~\ref{lem:linear-in-nu}, the map $\nu\mapsto \Lambda(\cdot\Vert\nu)$ is affine.
    Over the convex compact set $\cP(\cA\times\cL)$ an affine functional attains its supremum at an extreme point, i.e., a Dirac $\delta_{(a_\star,\ell)}$,
    and $\delta_{(a_\star,\ell)}\in\cN_{\mathrm{aff}}$ by the remark above.
    Hence the two suprema coincide.
\end{proof}

\begin{theorem}[Restricted-dual equality]\label{thm:restricted-dual}
    Fix $a,\eta$ and suppose Assumption~\ref{ass:compact} holds. Then
    \[
        \inf_{P\in\Delta(\cA),\,T\in\cT}\;\sup_{\nu\in\cP(\cA\times\cL)} \Lambda_{a,\eta}(P,E_T\Vert \nu)
        \;=\;
        \inf_{P\in\Delta(\cA),\,T\in\cT}\;\sup_{\nu\in\cN_{\mathrm{aff}}} \Lambda_{a,\eta}(P,E_T\Vert \nu).
    \]
    Consequently, combining with Lemma~\ref{lem:exp-by-opt-minimax} (Sion) gives
    \[
        \sup_{\nu\in\cP(\cA\times\cL)}\;\inf_{P\in\Delta(\cA),\,T\in\cT} \Lambda_{a,\eta}(P,E_T\Vert \nu)
        \;\le\;
        \inf_{P\in\Delta(\cA),\,T\in\cT}\;\sup_{\nu\in\cN_{\mathrm{aff}}} \Lambda_{a,\eta}(P,E_T\Vert \nu),
    \]
    i.e., the standard dual upper bound on the primal value is unchanged if the dual adversary is restricted to $\cN_{\mathrm{aff}}$.
\end{theorem}

\begin{proof}
    The first equality is immediate from Proposition~\ref{prop:sup-over-aff} by taking the infimum over $(P,T)$. The inequality is the standard minimax bound.
\end{proof}

\paragraph{What this buys us.}
Inside the RHS of Theorem~\ref{thm:restricted-dual} we may now \emph{pick} $\nu\in\cN_{\mathrm{aff}}$ and choose $a$ and $T$ so that $\tilde{A}=\po$ almost surely (Prop.~\ref{prop:mulY-po-iff}).
By Lemma~\ref{lem:E-optimized-Lambda} and Lemma~\ref{lem:the-right-turn},
\[
    \Lambda(P,E_T\Vert \nu)
    =
    \frac{1}{\eta}\Delta(\nu,P)
    -\frac{1}{\eta^2}\E\!\left[D_F(\po,a)\right]
    =
    \frac{1}{\eta}\Delta(\nu,P)
    -\frac{1}{\eta^2}\Big(\E[D_F(\po,\pr)]+D_F(\pr,a)\Big),
\]
whence optimizing $\eta$ yields the \emph{same} information-ratio upper bound
as in Theorem~\ref{thm:optimal-Lambda-info-ratio} for unrestricted estimators:
\[
    \Lambda(P,E_T\Vert \nu)\;\le\;\frac{\Delta(\nu,P)^2}{4\,\E[D_F(\po,\pr)]}.
\]
Therefore, \emph{for the dual analysis that upper-bounds regret}, multiplicative-$Y$ estimators are as powerful as unrestricted estimators.
\medskip

\paragraph{Roadblock: restricting the \emph{primal} (sup--inf) side.}
A natural wish is to also replace $\sup_{\nu\in\cP(\cA\times\cL)}$ by $\sup_{\nu\in\cN_{\mathrm{aff}}}$ \emph{before} swapping the order via Sion. This runs into two issues:

\begin{enumerate}
    \item \textbf{Non-convexity of $\cN_{\mathrm{aff}}$.} If $\nu_1,\nu_2\in\cN_{\mathrm{aff}}$ have different dual intercepts/slopes, their mixture need not remain dual-affine (e.g., mixtures of linear-Gaussian models are generally non-Gaussian and the posterior mean in $y$ becomes nonlinear), so Sion's convexity requirement for the adversary's set fails.

    \item \textbf{Concavity of the value map.} The map
          $f(\nu)\triangleq \inf_{P,T}\Lambda(P,E_T\Vert\nu)$ is the pointwise infimum of affine functionals in $\nu$, hence \emph{concave} (and upper semicontinuous under Assumption~\ref{ass:compact}). For concave $f$, $\sup$ over a nonconvex subset can be strictly \emph{smaller} than $\sup$ over its convex hull, so restricting the \emph{primal} $\sup$ to $\cN_{\mathrm{aff}}$ may strictly lower the value and break minimax equality.
\end{enumerate}

\paragraph{Two robust sub-classes where the affine property holds generically.}
Although $\cN_{\mathrm{aff}}$ is not convex in general, there are rich families where it is easy to operate:

\begin{itemize}
    \item \textbf{Quadratic $F$ (linear-Gaussian priors).} If $(A_\star,\ell)$ is jointly Gaussian (any mean and covariance) then for \emph{every} $b$,
          $(A_\star,Y_b)$ is jointly Gaussian and $\E[A_\star\mid Y_b=y]=U+S_b\,y$ is affine in $y$ with a common intercept $U$. Hence $\nu\in\cN_{\mathrm{aff}}$.

    \item \textbf{Negentropy (log-linear posteriors).} If for each $b$ the conditional law of $A_\star\mid Y_b=y$ is such that
          $\ln\E[A_\star\mid Y_b=y]$ is affine in $y$ (e.g., suitable exponential-family couplings in dual coordinates), then $\nu\in\cN_{\mathrm{aff}}$ and multiplicative-$Y$ recovers $\tilde{A}=\po$ a.s.\ via $T_b=-S_b/\eta$.
\end{itemize}

\paragraph{Takeaway.}
For regret \emph{upper bounds}, it suffices to analyze the dual problem and, by Theorem~\ref{thm:restricted-dual}, the adversary can be restricted to $\cN_{\mathrm{aff}}$ without changing the dual value. Inside this class, multiplicative-$Y$ achieves $\tilde{A}=\po$ a.s.\ and inherits the information-ratio bound of the unrestricted estimator. Attempting to impose the same restriction on the \emph{primal} side runs into genuine obstacles (non-convexity of priors with dual-affine posteriors and the concavity of $f(\nu)$).

\section{Nov 4, try 1}
\subsection{When can multiplicative-\texorpdfstring{$Y$}{Y} recover \texorpdfstring{$A_\star$}{A\_star} (or the posterior mean)?}

\paragraph{Setup.}
Throughout this subsection we work with the multiplicative-$Y$ class
$E_T(x,y)=y\,T(x)$, hence $\hl=Y\,T_A$ and
\begin{equation}
    \tilde{A} \;=\; {F^\star}'\!\big(F'(a)-\eta\,Y\,T_A\big).
    \label{eq:tildeA-mulY}
\end{equation}

\begin{lemma}[Measurability and a necessary condition]\label{lem:measurability}
    In the multiplicative-$Y$ class, $\tilde{A}$ is $\sigma(A,Y)$-measurable.
    Consequently, if $\tilde{A}=A_\star$ almost surely, then the posterior is
    degenerate: $\Var(A_\star\mid A,Y)=0$ almost surely (equivalently,
    $\po=A_\star$ a.s.).
\end{lemma}

\begin{proof}
    By \eqref{eq:tildeA-mulY}, $\tilde{A}$ is a measurable function of $(A,Y)$.
    If $\tilde{A}=A_\star$ a.s., then $A_\star$ is $\sigma(A,Y)$-measurable,
    which is equivalent to $\Var(A_\star\mid A,Y)=0$ a.s.
\end{proof}

\begin{proposition}[Characterization: $\tilde{A}=\po$ a.s. in the multiplicative-$Y$ class]
    \label{prop:mulY-po-iff}
    Assume $F$ is Legendre so that $F'$ is injective. Using $\hl=Y\,T_A$, we have
    \[
        \tilde{A}=\po \;\; \text{almost surely}
        \quad\Longleftrightarrow\quad
        F'\!\big(\po\big) \;=\; F'(a)\;-\;\eta\,Y\,T_A \;\; \text{almost surely.}
    \]
    Equivalently, for every $b$ with $P(A=b)>0$, the map
    $y\mapsto F'\!\big(\E[A_\star\mid A=b,Y=y]\big)$ is (a.s.) \emph{affine in $y$}
    with intercept $F'(a)$ and slope $-\eta\,T_b$.
\end{proposition}

\begin{proof}
    From \eqref{eq:tildeA-mulY}, $\tilde{A}=\po$ a.s.\ iff
    \(
    \po={F^\star}'\!\big(F'(a)-\eta\,Y\,T_A\big)
    \)
    a.s. Applying $F'$ (injective) gives the stated condition.
\end{proof}

\begin{remark}[Common intercept requirement]
    \label{rem:common-intercept}
    Since $a$ is fixed before sampling $A\sim P$, the “intercept”
    $F'(a)$ must be valid for all $b$ in $\mathrm{supp}(P)$
    in Proposition~\ref{prop:mulY-po-iff}. Thus, if for $b$ in a set
    $S\subseteq\cA$ the functions $y\mapsto F'(\po\mid A=b)$ share a common
    intercept $u$ (but may have different slopes), then it suffices to choose $a$
    so that $F'(a)=u$ and restrict $P$ to $S$.
\end{remark}

\begin{theorem}[Exact recovery $\tilde{A}=A_\star$ a.s.: dual-affine determinism]
    \label{thm:exact-recovery}
    Suppose there exists a vector $u$ and a function $T:\cA\to\R^d$ such that, for every
    $b$ in $\mathrm{supp}(P)$, the prior $\nu$ satisfies the \emph{dual-affine determinism}
    \[
        F'\!\big(A_\star\big) \;=\; u\;-\;\eta\,Y\,T_b
        \qquad\text{almost surely given } A=b.
    \]
    Choose $a$ with $F'(a)=u$ and use the multiplicative-$Y$ estimator with this $T$.
    Then $\tilde{A}=A_\star$ almost surely. In particular, $\po=A_\star$ a.s.\ and
    $\Var(A_\star\mid A,Y)=0$ a.s.
\end{theorem}

\begin{proof}
    By the assumed relation and \eqref{eq:tildeA-mulY},
    \(
    \tilde{A}={F^\star}'(u-\eta\,Y\,T_A)=A_\star
    \)
    a.s. The rest follows from Lemma~\ref{lem:measurability}.
\end{proof}

\begin{corollary}[Information-ratio bound for multiplicative-$Y$]
    \label{cor:IR-mulY}
    If $\tilde{A}=\po$ almost surely (e.g., by Proposition~\ref{prop:mulY-po-iff}
    or Theorem~\ref{thm:exact-recovery}), then
    \[
        \Lambda(P, E_T \,\Vert\, \nu)
        \;=\;
        \frac{1}{\eta}\,\Delta(\nu,P)
        \;-\;\frac{1}{\eta^2}\,\E\!\left[D_F\!\left(\po,a\right)\right],
    \]
    and hence the same information-ratio bound as in the unrestricted case holds:
    \[
        \Lambda(P, E_T \,\Vert\, \nu)
        \;\le\;
        \frac{\Delta(\nu,P)^2}{4\,\E\!\left[D_F(\po,\pr)\right]}.
    \]
\end{corollary}

\begin{proof}
    From Lemma~\ref{lem:E-optimized-Lambda} we have
    \(
    \Lambda(P,E_T\Vert\nu)
    = \frac{1}{\eta}\Delta(\nu,P)-\frac{1}{\eta^2}\E[D_F(\tilde{A},a)].
    \)
    If $\tilde{A}=\po$ a.s., then the expected Bregman identity yields
    \(
    \E[D_F(\po,a)]=\E[D_F(\po,\pr)]+D_F(\pr,a).
    \)
    Maximizing in $\eta$ gives the stated bound, identical to
    Theorem~\ref{thm:optimal-Lambda-info-ratio}.
\end{proof}

\paragraph{Two canonical templates for \texorpdfstring{$\nu$}{nu} where \texorpdfstring{$\tilde{A}=A_\star$}{A\_tilde=A\_star} is feasible.}

\begin{example}[Quadratic potential, single-index linear model]\label{ex:quad-single-index}
    Let $F(x)=\tfrac12\|x\|^2$ so $F'(x)=x$ and ${F^\star}'(q)=q$.
    Assume there exists a scalar latent $\theta$ and action-dependent slope $t:\cA\to\R^d$
    and a fixed $u\in\R^d$ such that, conditional on $A=b$,
    \[
        \ell = \theta\,s \quad\text{and}\quad
        A_\star = u - \gamma\,\theta\,t_b,
    \]
    for some fixed $s\in\R^d$ and $\gamma>0$.
    Then $Y=\langle b,\ell\rangle=\theta\,\langle b,s\rangle$ and
    \(
    A_\star = u - \eta\,Y\,\Big(\frac{\gamma}{\eta\,\langle b,s\rangle}t_b\Big).
    \)
    If we pick $a=u$ and $T_b=\frac{\gamma}{\eta\,\langle b,s\rangle}t_b$, then
    \(
    \tilde{A}=a-\eta Y T_A = u-\gamma\,\theta\,t_A=A_\star
    \)
    a.s.  By Cor.~\ref{cor:IR-mulY}, the information-ratio bound holds.
\end{example}

\begin{example}[Negentropy on the simplex, dual-log-linear posterior]\label{ex:negentropy-loglinear}
    Let $\cA=\Delta(d)$ and $F(p)=\langle p,\ln p -1\rangle$ so $F'(p)=\ln p$ and
    ${F^\star}'(q)=\exp(q)$ (coordinatewise; cf.\ Sec.~\ref{sec:md-negentropy}).
    Assume that for each $b$ in $\mathrm{supp}(P)$ there exists $t_b\in\R^d$ and a common
    $u\in\R^d$ such that, conditional on $A=b$,
    \[
        \ln A_\star \;=\; u \;-\;\eta\,Y\,t_b \qquad\text{almost surely.}
    \]
    Choosing $a$ with $\ln a=u$ and the multiplicative-$Y$ estimator with $T_b=t_b$ gives
    \(
    \tilde{A}=\exp(\ln a - \eta Y t_b)=a\odot \exp(-\eta Y t_b)=A_\star
    \)
    a.s.  As in Example~\ref{ex:quad-single-index}, Cor.~\ref{cor:IR-mulY} applies.
\end{example}

\paragraph{Impossibility outside the dual-affine family.}
Even for $\tilde{A}=\po$ (let alone $\tilde{A}=A_\star$), Proposition~\ref{prop:mulY-po-iff}
shows we need that, for each $b$ in $\mathrm{supp}(P)$,
$y\mapsto F'\!\big(\E[A_\star\mid A=b,Y=y]\big)$ is affine in $y$ with a \emph{common intercept}
$F'(a)$. If this fails on a set of positive probability (e.g., the posterior mean depends on $y$
nonlinearly in dual coordinates or the intercepts differ across $b$ in the support of $P$),
then no multiplicative-$Y$ estimator can make $\tilde{A}=\po$ a.s., and the clean
information-ratio reduction need not hold without additional assumptions
(cf.\ the counterexample in Thm.~\ref{thm:no-alpha}).

\begin{remark}[Practical prescription]
    To make the multiplicative-$Y$ class inherit the information-ratio bound, it suffices to:
    (i) restrict $P$ to a support $S$ on which the dual posterior means share a common intercept $u$,
    (ii) set $a$ so that $F'(a)=u$, and (iii) choose $T_b$ as the (negative) dual slope
    of $y\mapsto F'(\po\mid A=b,Y=y)$ for each $b\in S$. Under this “dual-affine posterior”
    condition, $\tilde{A}=\po$ a.s.\ and Cor.~\ref{cor:IR-mulY} applies.
\end{remark}

\hrule

\subsection*{When does \texorpdfstring{$\dA=\po$}{A\_tilde=po} almost surely?}

Recall the definition
\[
    \dA \;\triangleq\; {F^\star}'\!\left(F'(a)-\eta\,\hl\right),
    \qquad
    \po \;\triangleq\; \E\!\left[A_\star \,\middle|\, A, Y\right],
\]
and (in the multiplicative-$Y$ class) $\hl=\frac{Y}{\eta P_A}T_A$ with $T\in(\R^d)^{\cA}$.

\begin{corollary}[Quadratic potential $\Rightarrow$ linear regression in $Y$]
    \label{cor:quad-linear}
    Let $F(x)=\tfrac12\|x\|^2$. Then $\dA=\po$ a.s.\ iff, for $\PP$-a.e.\ $b$,
    \begin{equation}
        \label{eq:lin-reg}
        \E[A_\star\mid A=b,Y=y] \;=\; a - \gamma_b\,y
        \quad\text{for $Y$-a.e.\ $y$,}
    \end{equation}
    with some $\gamma_b\in\R^{d}$, in which case choosing $T_b=P_b\,\gamma_b$ yields $\dA=\po$ a.s.
    A canonical sufficient condition for \eqref{eq:lin-reg} is joint conditional Gaussianity:
    if $(A_\star,Y)\mid(A=b)$ is jointly Gaussian with
    \(
    \E[A_\star\mid A=b]=\pr,\;
    \E[Y\mid A=b]=\mu_b,
    \)
    and
    \(
    \Gamma_b=\Cov(A_\star,Y\mid A=b),\;
    \sigma_b^2=\Var(Y\mid A=b)>0,
    \)
    then
    \[
        \po \;=\; \pr + \Gamma_b\,\sigma_b^{-2}\,(Y-\mu_b).
    \]
    Setting
    \(
    T_b \;=\; -P_b\,\Gamma_b\,\sigma_b^{-2}
    \)
    and
    \(
    a \;=\; \pr - \Gamma_b\,\sigma_b^{-2}\,\mu_b
    \)
    forces \eqref{eq:lin-reg} and hence $\dA=\po$ a.s.
    In particular, if $\E[Y\mid A=b]=0$ (e.g.\ $\E[\ell]=0$ under $\nu$), it suffices to take $a=\pr$ and $T_b=-P_b\,\Gamma_b\,\sigma_b^{-2}$.
\end{corollary}

\begin{corollary}[Mirror–exponential-family conjugacy]
    \label{cor:mirror-conjugacy}
    Assume $F$ is Legendre and consider the exponential family on $\cA$ with log-partition $F^\star$:
    \(
    \pi_\theta(\d x)\propto h(x)\,\exp\{\langle\theta,x\rangle - F^\star(\theta)\}\,\d x.
    \)
    Suppose $A_\star\sim \pi_{\theta_0}$ with $\theta_0=F'(a)$, and, for each $b$,
    the conditional law of $Y$ given $(A_\star=x,A=b)$ has density
    \[
        q_b(y\mid x) \;\propto\; r_b(y)\,\exp\!\left\{-\tfrac{y}{P_b}\,\langle T_b, x\rangle\right\},
    \]
    for some measurable $r_b$ and fixed $T_b\in\R^d$.
    Then Bayes' rule gives the posterior
    \(
    A_\star\mid (A=b,Y=y)\sim \pi_{\theta_0-(y/P_b)T_b},
    \)
    hence
    \(
    \po={F^\star}'\!\big(\theta_0-(y/P_b)T_b\big)
    ={F^\star}'\!\left(F'(a)-\tfrac{y}{P_b}T_b\right)=\dA
    \)
    a.s.
\end{corollary}

\begin{corollary}[Degenerate $Y$]
    \label{cor:degenerate}
    If $Y$ is (conditionally) degenerate given $A=b$ (i.e.\ $\PP(Y=y_b\mid A=b)=1$), then there exists $T_b$ with
    \(
    F'\!\big(\po\big)=F'(a)-\tfrac{y_b}{P_b}T_b
    \)
    and consequently $\dA=\po$ a.s.\ (this is vacuous for learning, but it shows attainability in degenerate cases).
\end{corollary}

\paragraph{A necessary condition (no hidden curvature in \texorpdfstring{$y$}{y}).}
\begin{lemma}[Affine-dual necessity]
    \label{lem:affine-necessary}
    Suppose $E_T$ is multiplicative-$Y$ and $\dA=\po$ a.s.
    Fix $b$ with $\PP(A=b)>0$ and assume $y\mapsto \po(b,y)$ is a.s.\ differentiable.
    Then the \emph{dual} posterior mean must be affine in $y$:
    \[
        \frac{\partial^2}{\partial y^2}\,F'\!\big(\po(b,y)\big) \;=\; 0
        \quad\text{for $Y$-a.e.\ $y$.}
    \]
    For $F(x)=\tfrac12\|x\|^2$, this reduces to the linear-regression property
    $\frac{\partial^2}{\partial y^2}\E[A_\star\mid A=b,Y=y]=0$ a.s.
\end{lemma}

\paragraph{Why the condition is restrictive.}
In general, the map $y\mapsto \po$ need \emph{not} be affine (in primal or dual coordinates).%
\footnote{See the counterexample constructed at the end of the document, where the posterior mean contains a quadratic term in the latent $Z$; there the ratio $\E[D_F(\dA,\pr)]/\E[D_F(\po,\pr)]$ can be driven to $0$, showing that even a uniform multiplicative bound fails in general.}
This explains why the equality $\dA=\po$ cannot be enforced by multiplicative-$Y$ estimators unless one of the structural scenarios above holds.

\paragraph{Summary.}
\begin{itemize}
    \item With the \emph{unrestricted} optimal estimator $G_{\nu,P}$, one always has $\dA=\po$ a.s.
    \item Within the \emph{multiplicative-$Y$} class, $\dA=\po$ a.s.\ \emph{iff} the conditional posterior mean evolves along a straight line in \emph{dual} space as $y$ varies (Prop.~\ref{prop:multY-iff}). This is guaranteed by (i) quadratic $F$ $+$ linear conditional regression (e.g.\ conditional Gaussianity) or (ii) mirror–exponential-family conjugacy where observing $Y$ induces an additive shift of the natural parameter.
    \item Absent such structure, $\dA=\po$ a.s.\ typically fails.
\end{itemize}


\hrule

\subsubsection{Optimizing for \texorpdfstring{$E$}{E}---linear (affine) in \texorpdfstring{$Y$}{Y} without normalizers}
We now remove the explicit normalizers from the estimator and work directly with
\begin{equation}
    \label{eq:def-linY-no-norm}
    E_{T,S}(a,y)
    \;=\;
    T_a\,y \;+\; S_a,
    \qquad
    \hl \;=\; E_{T,S}(A,Y),
\end{equation}
for matrices $T,S \in \R^{|\cA|\times d}$.  As before, define
$\tilde{A}\triangleq {F^\star}'\!\big(F'(a)-\eta \hl\big)$ and use
$\Lambda(P,E\Vert \nu)$ from \eqref{eq:Lambda-simplified}.  This parameterization
is equivalent to the previous one (with $1/(\eta P_a)$) in the sense of
Remark~\ref{rem:linY-absorb} below, but makes the presentation cleaner.

\paragraph{First-order optimality with slope \& intercept.}
For fixed $a,\eta,P,\nu$, consider
\(
\inf_{T,S} \Lambda(P, E_{T,S}\Vert \nu)
\).
The optimality conditions mirror those in the multiplicative class.

\begin{theorem}[Linear-$Y$ optimality conditions (no normalizers)]
    \label{thm:optimal-linY-no-norm}
    Let $(T^\star,S^\star)$ minimize $\Lambda(P,E_{T,S}\Vert \nu)$ over the affine
    class \eqref{eq:def-linY-no-norm} and set $\hl=E_{T^\star,S^\star}(A,Y)$ and
    $\tilde{A}={F^\star}'(F'(a)-\eta \hl)$. Then for every $b\in\cA$,
    \begin{align}
        \E\!\left[\, Y\,(A_\star-\tilde{A}) \,\big|\, A=b \right] & = 0,
        \label{eq:linY-FOC-slope-no-norm}
        \\
        \E\!\left[\, A_\star-\tilde{A} \,\big|\, A=b \right]      & = 0.
        \label{eq:linY-FOC-intercept-no-norm}
    \end{align}
    Equivalently, conditional on $A=b$, the residual $A_\star-\tilde{A}$ is
    orthogonal to both the constant function $1$ and to the scalar $Y$.
\end{theorem}

\begin{proof}
    Starting from \eqref{eq:Lambda-simplified} and conditioning on $A=b$,
    minimize
    \[
        \Phi_b(T_b,S_b)
        =
        \E\!\left[
            \left\langle A_\star, T_b Y + S_b \right\rangle
            + \tfrac{1}{\eta} F^\star\!\left(F'(a)-\eta(T_b Y + S_b)\right)
            \,\Big|\, A=b
            \right].
    \]
    Differentiating and setting the gradients to zero yields
    \[
        0=\tfrac{\partial \Phi_b}{\partial T_b}
        \;\Rightarrow\;
        \E\!\left[ Y A_\star - Y\,{F^\star}'\!\left(F'(a)-\eta(T_b Y + S_b)\right)\Big|A=b\right]=0,
    \]
    and
    \[
        0=\tfrac{\partial \Phi_b}{\partial S_b}
        \;\Rightarrow\;
        \E\!\left[ A_\star - {F^\star}'\!\left(F'(a)-\eta(T_b Y + S_b)\right)\Big|A=b\right]=0.
    \]
    With $\hl=T_b Y + S_b$ these are exactly
    \eqref{eq:linY-FOC-slope-no-norm}--\eqref{eq:linY-FOC-intercept-no-norm}.
\end{proof}

\begin{lemma}[Calibration and replacement identity]
    \label{lem:linY-calibration-replace-no-norm}
    Under the conditions of Theorem~\ref{thm:optimal-linY-no-norm}:
    \begin{enumerate}
        \item (\emph{Calibration}) Since $A$ is independent of $(A_\star,\ell)$ under $\nu$,
              \eqref{eq:linY-FOC-intercept-no-norm} implies
              \[
                  \E[\tilde{A}\mid A=b] \;=\; \E[A_\star\mid A=b] \;=\; \pr,\qquad\forall b\in\cA.
              \]
        \item (\emph{Replacing $A_\star$ by $\tilde{A}$ inside $\langle\cdot,\hl\rangle$})
              Using \eqref{eq:linY-FOC-slope-no-norm}--\eqref{eq:linY-FOC-intercept-no-norm} and
              $\hl=T_A Y + S_A$,
              \[
                  \E\!\left[\langle A_\star,\hl\rangle\mid A\right]
                  \;=\;
                  \E\!\left[\langle \tilde{A},\hl\rangle\mid A\right].
              \]
    \end{enumerate}
\end{lemma}

\begin{lemma}[E-optimized $\Lambda$ (linear-$Y$, no normalizers)]
    \label{lem:linY-optimized-Lambda-no-norm}
    Let $(T^\star,S^\star)$ be optimal and $\hl=E_{T^\star,S^\star}(A,Y)$. Then
    \begin{equation}
        \Lambda(P,E_{T^\star,S^\star}\Vert \nu)
        \;=\;
        \frac{1}{\eta}\,\Delta(\nu,P)
        \;-\;
        \frac{1}{\eta^2}\,
        \E\!\left[
            D_{F^\star}\!\left(F'(a),\,F'(a)-\eta \hl\right)
            \right].
        \label{eq:linY-optimized-Lambda-no-norm}
    \end{equation}
\end{lemma}

\begin{proof}
    From \eqref{eq:Lambda-simplified} and
    Lemma~\ref{lem:linY-calibration-replace-no-norm}(ii),
    replace $\langle A_\star,\hl\rangle$ by $\langle \tilde{A},\hl\rangle$ and use
    the identity
    \[
        \big\langle {F^\star}'(F'(a)-\eta\hl),\hl\big\rangle
        +\tfrac{1}{\eta}F^\star(F'(a)-\eta\hl) - \tfrac{1}{\eta}F^\star(F'(a))
        = -\,\tfrac{1}{\eta}\,D_{F^\star}(F'(a),F'(a)-\eta\hl).
    \]
\end{proof}

\paragraph{Orthogonality and the expected-Bregman identity (unchanged).}
Using $F'(a)-F'(\tilde{A})=\eta \hl$ and
Lemma~\ref{lem:linY-calibration-replace-no-norm}, we get
\[
    \E\!\left[
        \big\langle F'(a)-F'(\tilde{A}),\, \po-\tilde{A}\big\rangle
        \,\Big|\, A
        \right]
    = \eta\,\E\!\left[\langle \hl,\, \po-\tilde{A}\rangle \big| A\right]
    = 0.
\]
Applying the three-point identity and averaging over $A$ recovers exactly
\eqref{eq:expected-bregman-identity}:
\begin{equation}
    \E\!\left[D_F(\po,\tilde{A}) + D_F(\tilde{A},a)\right]
    \;=\;
    \E\!\left[D_F(\po,\pr)\right] + D_F(\pr,a).
    \label{eq:expected-bregman-identity-linY-no-norm}
\end{equation}
Hence,
\[
    \E[D_F(\tilde{A},a)]
    \;=\;
    \E[D_F(\po,\pr)] - \E[D_F(\po,\tilde{A})] + D_F(\pr,a),
\]
and the same discussion as in the multiplicative case applies: without further
structure, $\E[D_F(\tilde{A},a)]$ need not lower bound
$\E[D_F(\po,\pr)]$ by a universal multiplicative factor.

\begin{corollary}[No universal multiplicative factor, even without normalizers]
    \label{cor:no-alpha-linY-no-norm}
    The counterexample under \emph{Multiplicative factor?} continues to apply to the
    affine class \eqref{eq:def-linY-no-norm}.  In particular, with quadratic $F$
    and $\E[Y\mid A]=0$, the intercept condition
    \eqref{eq:linY-FOC-intercept-no-norm} calibrates the mean, while the slope is
    fixed by \eqref{eq:linY-FOC-slope-no-norm}; one can still make
    $\E[D_F(\tilde{A},\pr)]/\E[D_F(\po,\pr)]\to 0$, so no universal $\alpha>0$
    exists.
\end{corollary}

\paragraph{When can we achieve \texorpdfstring{$\tilde{A}=\po$}{A\_tilde=po} a.s.?}
Allowing a bias relaxes the intercept but does not remove the fundamental
nonlinearity barrier.

\begin{proposition}[Characterization in the linear-$Y$ class (no normalizers)]
    \label{prop:linY-iff-no-norm}
    Consider $E_{T,S}$ with $\hl=T_A Y + S_A$ and
    $\tilde{A}={F^\star}'(F'(a)-\eta \hl)$. Then $\tilde{A}=\po$ a.s.\ if and only if
    \begin{equation}
        \label{eq:affine-dual-linY-no-norm}
        F'\!\big(\po\big)
        \;=\;
        F'(a) \;-\; \eta\big(T_A\,Y + S_A\big)
        \qquad\text{almost surely.}
    \end{equation}
    Equivalently, conditional on $A=b$, the map
    $y\mapsto F'\!\big(\E[A_\star\mid A=b,Y=y]\big)$ must be \emph{affine in $y$},
    with slope $-\eta T_b$ and intercept $F'(a)-\eta S_b$.
\end{proposition}

\begin{proof}
    $\tilde{A}=\po$ a.s.\ $\iff$
    $\po={F^\star}'\!\big(F'(a)-\eta \hl\big)$ a.s.  Apply $F'$ and substitute
    $\hl=T_A Y + S_A$.
\end{proof}

\begin{remark}[Absorbing normalizers and parameterization invariance]
    \label{rem:linY-absorb}
    The class $\hl=T_A Y+S_A$ is equivalent to the earlier
    $\hl=\frac{Y}{\eta P_A}\,\widetilde{T}_A+\widetilde{S}_A$:
    for any $(\widetilde{T},\widetilde{S})$ define
    $T_b=\widetilde{T}_b/(\eta P_b)$ and $S_b=\widetilde{S}_b$ (for $P_b>0$) to get
    the same $\hl$; conversely $\widetilde{T}_b=\eta P_b T_b$,
    $\widetilde{S}_b=S_b$.  Therefore the achievable values of
    $\Lambda(P,E\Vert \nu)$ and all bounds obtained via
    \eqref{eq:linY-optimized-Lambda-no-norm} are unchanged by this reparameterization.
\end{remark}

\begin{remark}[Quadratic potential sanity check]
    For $F(x)=\frac12\|x\|^2$, we have $F'(x)=x$ and
    $\tilde{A}=a-\eta(T_A Y+S_A)$.  If additionally $\E[Y\mid A]=0$, the
    intercept condition forces $S_b=(a-\pr)/\eta$ and thus
    $\tilde{A}=-\eta\,T_A Y$; i.e., the intercept only calibrates the mean and the
    slope $T$ drives the informational term in
    \eqref{eq:linY-optimized-Lambda-no-norm}.
\end{remark}
